%%%%%%%%%%%%%%%%%%%%%%%%%%%%%%%%%%%%%%%%%%%%%%%%%%%%%%%%%%%%%%%%%%%%%%
% Lớp 10 | Bất phương trình bậc 2 một ẩn | Phương trình quy về phương trình bậc hai
% Dạng 0D7?3-3: Phương trình căn thức có tham số
% Số câu: 2  (sinh tự động bởi tach_cau.py; từ nay thầy sửa trực tiếp file này)
%%%%%%%%%%%%%%%%%%%%%%%%%%%%%%%%%%%%%%%%%%%%%%%%%%%%%%%%%%%%%%%%%%%%%%

%%%%%%%%%%%%%%%% TRẢ LỜI NGẮN %%%%%%%%%%%%%%%%

% nguồn: MaHoaNH_TraLoiNgan_Lop10
\begin{bt}%[0D7V3-3]%[ma:794f4560]
	Cho bất  phương trình $4\sqrt{(x+1)(3-x)}\le x^2-2x+m-3$. Có bao nhiêu giá trị nguyên của tham số $m\in (-2022;2023]$ để bất phương trình nghiệm đúng với $\forall x\in [-1;3]$?
	 \shortans{2012}
	 \loigiai{
		Với mọi $x\in [-1;3]$ ta  có $(x+1)(3-x)=-x^2+2x+3\ge 0$.\\
		Đặt $t=\sqrt{(x+1)(3-x)} \ge 0$, đồng  thời $\sqrt{(x+1)(3-x)}\le\dfrac{x+1+3-x}{2}=2$. Do đó, $t\in[0;2]$.\\
		Khi  đó bất phương trình trở thành $t^2+4t-m\le 0\Leftrightarrow m\ge t^2+4t$ với $t\in [0;2]$.\\
		Với $t\in [0;2]$ ta có  $0\le t^2+4t\le 12$.\\
		Do đó, $m\ge  12$.\\
		Vậy số giá trị nguyên của tham số $m$ là $2023-12+1=2012$. 
	 }
\end{bt}

% nguồn: MaHoaNH_TraLoiNgan_Lop10
\begin{bt}%[0D7C3-3]%[ma:0ff0285f]
	Tìm số giá trị nguyên của  tham số $m$ thuộc đoạn $[0;2024]$ thoả mãn phương trình $\sqrt{2x^2 - 2x - 2m} = x-2$ có nghiệm.
	 \shortans{2023}
	\loigiai{ 
		Ta  có
		\[ 
		\sqrt{2x^2 - 2x -  2m} = x-2 \Leftrightarrow \heva{&x-2\ge 0\\ &2x^2 - 2x - 2m = (x-2)^2}
		 \Leftrightarrow \heva{&x\ge 2\\ &x^2+2x-4 = 2m.}
		\] 
		Do đó  phương trình $\sqrt{2x^2 - 2x - 2m} = x-2$ có nghiệm khi và chỉ khi phương trình $x^2+2x-4 = 2m$ có nghiệm $x_0\ge 2$.\\
		Xét tam thức bậc hai $f(x) = x^2+2x-4$ trên $\left[2;+\infty \right)$ có $a = 1>0$ và hoành độ đỉnh  $x = -\dfrac{b}{2a} =-1 < 2$ nên ta có
		\[ 
		f(x) \ge f(2) =  4,\, \forall x\ge 2.
		 \]
		Suy  ra $2m\ge 4 \Leftrightarrow m\ge 2$.\\
		Kết hợp với $m\in[0;2024]$ ta thấy có  $2023$ giá trị nguyên của $m$ thoả mãn.
	 }
\end{bt}
